
\chapter{Literature Review}




\subsection{Current e-Governance & Complaint Systems}
\subsection{Blockchain in Governance}
\subsection{Identity Management}

% \subsection{Blockchain-Based Citizen Services and E-Governance}

% The adoption of blockchain  in e-governance has moved beyond theoretical frameworks to practical applications in citizen services. Case studies reveal implementations in feedback and voting systems, vehicle registration, e-tendering, complaint management and other paperless government services. For example DIGIT suite developed by the eGov Foundation serves as a comprehensive open-source platform for services as property tax management and birth/death registration, focusing on interoperability through open APIs. During the COVID-19 pandemic, blockchain was even used to crowdsource neighborhood status reports, demonstrating the technology's ability to achieve traceability and immutability in emergency public management. 

% \subsection{Decentralised Identity and \ac{SSI}}

% A fundamental requirement for digital governance is a secure identity system. \ac{SSI} is a decentralized framework that prevents central authorities from controlling identity data, instead of giving users full control. This system is typically anchored on a distributed ledger, allowing users to share only the minimal information required for a transaction. Research by Naik et al. has established twenty governing principles to evaluate the effectiveness and privacy-preserving capabilities of these \ac{SSI} ecosystems. 

% \subsection{Trust and Moderation Mechanisms}

% Existing literature has explored trust-level-based blockchain applications in specilized domains like the \ac{IoV} and social media news verification. these systems assign trust values to users based on the accuracy of their input compared to majority decisions. Furthermore, for managing crowdsourced content, there is a research into integrating \ac{AI} tools, such as Amazon Rekognition or openAI moderations, to filter harsh language or unsolicited imagery without slowing down the reporting process. 


% Different case studies of blockchain-based e-Governance
% services are discussed in \cite{9587836} and \cite{hamill2018blockchain}. They cover different citizen services as well as regulatory, compliance, and business operations. Citizen services extend to areas such as feedback and voting systems \cite{BAUDIER2021120397}, vehicle registration \cite{10338883}, e-tendering, complaint management\cite{9315884}, and other paperless government services. Platform to crowdfund municipal bonds, distributed energy markets, administering stocks, and track and traced controlled substances such as Cannabis fals under regulatory, compliance and business operations. Some of these services are fully implemented in countries such as Estonia
% and some are in pilot stage in different part of the world.
% These services take inherent advantage such as immutability
% and transparency of the underlying blockchain technology.
% However, some of the critical aspects such as user anonymity
% required in reporting services are overlooked in the police
% complaint management system. Trust levels of participants,
% security of the platforms, user privacy, existing legislation,
% public attitudes, operational cost, and technology penetration
% are identified as major barriers to implement such systems.
% A comprehensive open-source e-Governance suite, DIGIT,
% is implemented by the eGov Foundation \cite{egov_digit}. It provides
% a wide variety of tools for citizen-centric services such as
% birth and death registration and property tax management and
% local government administration. Further, it pays attention to
% core platform capabilities such as authentication and interoperability using open APIs. Although DIGIT delivers significant
% advancements in the space of digital public infrastructure, the
% documentation do not explicitly address the immutability of
% records or issue of preserving anonymity of users, which are
% paramount for securing citizen identities and fostering trust in
% sensitive reporting scenarios.


% \noindent A survey on existing citizen complaint management systems
% are presented in \cite{10566450} , which provides an insight to classical
% systems and their evolution to sophisticated platforms. Further,
% they discuss challenges for enhancing transparency and efficiency
% in government-citizen interactions and how emerging
% technologies such as \ac{AI} can be used to
% improve the service. Requirements of our target system such
% as traceability, immutability, and user privacy are not discussed
% in detail.

% \noindent Crowdsourced community governance was attempted during
% the COVID pandemic by encouraging citizens to report the status of their neighborhood\cite{9342128}. Participants are rewarded
% to sustain the service leading to successful identification of
% COVID spread within the area. Although immutability and
% traceability is achieved via blockchain, ownership of the
% blockchain by a particular stakeholder may raise concerns
% regarding the data integrity. Furthermore, the privacy of the
% reporting citizens or the data hasn’t been a primary concern.


% \noindent The \ac{SSI} is a decentralized identity
% framework which prevents the risks of identity being
% controlled by a central authority while giving the user the
% full control over their identity data\cite{9858139}. It is anchored by
% a distributed ledger technology such as blockchain and allows
% uses to share minimal amount of information required for
% authenticating a particular transaction. Naik et. al \cite{9272212} proposed
% twenty governing principles for SSI to analyze any SSI
% ecosystem and its effectiveness.


\section{Gaps in Literature}
\subsection{The Lack of User Anonymity and Control in Crowdsourced Governance}
\subsection{Overstated Data Integrity and Unaddressed Security Vulnerabilities}
\subsection{Lack of Trustworthiness Assessment for Individual Participants}
\subsection{Scalability and Operational Barriers in Developing Regions}
% \subsection{Neglect of User Anonymity and Protection}

% A significant gap exists regarding user anonymity in e-governance reporting services. While platforms like DIGIT and various police complaint management systems ensure data authentication and transparency, they often overlook the risk of retaliation against reporters. Without explicit anonymity, citizens reporting local issues remain vulnerable to the same targeting faced by whistleblowers. Currently technological surveys on complaint management systems often fail to discuss user privacy and traceability in suffiecient detail. 

% \subsection{Data Integrity and Stakeholder Centralization}

% While blockchain provide immutability, the ownership of the network presents a critical gap. In some community governance attempts, the blockchain was owned by a single stakeholder, which raises serious concerns regarding data integrity and the potential for manipulation. There is a lack of frameworks that balance the benefits of blockchain with a multi-stakeholder \ac{PoA} model involving independent agencies like the United Nations or Transparency International to ensure neutral validation. 

% \subsection{Economic and Technical Barriers in Developing Nations}

% Existing literature often ignores the financial constraints of implementing blockchain in **developing countries**. High **transaction (gas) costs** on public networks like Ethereum make them unsuitable for local governance in resource-constrained environments,. Furthermore, many current smart contract designs are not optimized, requiring "infinite" gas for execution. There is a clear need for research into **permissioned blockchains** that use fixed maintenance costs instead of per-transaction fees to make these systems viable for local administrations.

% \subsection{Reliability of Crowdsourced Input}

% The reliability of community-assisted services is currently dictated by the integrity of the individuals involved, yet there is a lack of robust mechanisms to handle **users with vested interests**. While trust-level applications exist in other fields, they have not been widely integrated into local governance reporting to filter out biased or fraudulent reports from the public.
